\chapter{Concluzii și perspective}
\label{chapter:concluzii}

\section{Concluzii}
\label{sec:concluzii}

În concluzie, proiectul VisioScience3D reprezintă o soluție completă pentru gestionarea și monitorizarea performanțelor 
elevilor la materiile de științe sau în termeni internaționali STEM. Platforma integrează funcționalități robuste care
asigură o experiență interactivă și captivantă utilizatorilor, demonstrând o retenție puternică în utilizare. Soluția e
bazată pe o infrastructură cu microservicii construită pe Kubernetes. Ea expune o interfață interactivă și un back-end
scalabil după cum am descris și oferă astfel o soluție de viitor pentru nevoile elevilor și profesorilor.

Testele realizate au demonstrat că soluția implementată îndeplinește obiectivele stabilite,
iar performanța și stabilitatea soluției trec evaluările de integrare finale făcute pe majoritatea scenariilor
cu utilizatori reali. De asemenea, expunerea noțiunilor prin componente 3D și a funcționalități interactive a
aduce cu siguranță un plus de valoare aplicației. Feedback-ul per total dat de testeri a fost foarte pozitiv în
ceea ce privește ușurința în utilizare și performanța sistemului.

\section{Dezvoltare viitoare}
\label{sec:dezvoltare-viitoare}


Pentru a îmbunătăți în viitor aplicația și a răspunde unor nevoi mai complexe, există mai multe direcții
posibile de dezvoltare viitoare:

Funcționalități de export și analize suplimentare pe care le pot face profesorii:
Ar putea fi adăugat un buton de export a rezultatelor din leaderboard-uri și de la teste în general.
S-ar putea adăuga suport pentru exportul datelor folosind formate precum CSV, PDF sau altele la cerere.

Integrarea de notificări și alerte: 
Implementarea unor notificări prin email pentru utilizatori ar putea îmbunătăți
experiența acestora, anunțând despre teste viitoare, termene limită sau evaluări de performanță.

Integrarea cu modele populare de inteligență artificială (precum modele lingvistice de exemplu):
Implementarea unor funcționalități de suport pentru elevi bazat pe inteligență artificială
ar putea sprijini foarte mult procesul de învățare, oferindu explicații detaliate și sugestii personalizate
pentru întrebările elevilor de exemplu.

Funcționalități de colaborare:
Oferirea unui sistem de chat în timp real integrat în aplicație ar putea facilita colaborarea între utilizatori, 
eliminând necesitatea utilizării unor alte platforme externe pentru comunicare.

Prin urmare, viitorul aplicației se axează pe extinderea funcționalităților, îmbunătățirea performanței și
scalabilității și integrarea de noi tehnologii pentru a oferi utilizatorilor o experiență din ce în ce
mai personalizată și interactivă.